\subsection*{Drill: Three Wall Double Lines}
\label{drill:three_wall/double_lines}

\begin{minipage}{0.6\textwidth}

This drill involves a three wall and a jammer.

There are two variations of this drill.
In rough order of difficulty these are:
\begin{itemize}
\item Blockers in lanes 1, 2, brace between lanes 1 and 2 
\item Blockers in lanes 3, 4, brace between lanes 3 and 4 
\end{itemize} 

The jammer begins in the seam. 

The goal of the jammer is to out-pace the wall to the other side of the track and pass the blockers in either lane 1 or lane 4.
The jammer must commit to this movement until they have passed the first blocker. 

The jammer's second movement should then be to cut back in. 


\end{minipage}
\hfill
\hspace{0.5cm}
\begin{minipage}{0.38\textwidth}
\begin{centering}
\includegraphics[scale=1.1,trim={0 0 0 4cm}, clip]{img/double_lines/escape}
{\it Jammer cuts back in, avoiding the trap on the line.} 
\end{centering}
\end{minipage}

\vspace{0.5cm}
\begin{minipage}{0.38\textwidth}
\begin{centering}

\includegraphics[scale=1.1,trim={0 0 0 4cm}, clip]{img/double_lines/caught}
{\it Jammer takes a double line and is caught or hit out.} 
\end{centering}
\end{minipage}
\hfill
\hspace{0.5cm}
\begin{minipage}{0.6\textwidth}

The goals of this drill vary depending on the level of the jammer and the blockers.

\begin{itemize}
\item The first goal is for the jammer to outpace the blockers sufficiently to escape.
\item The second is for the blockers to reach the line in time to prevent this. This will often involve the first blocker arriving too late, while the brace is in position to catch the jammer.
\item The third goal is for the jammer to avoid taking a double line and instead cut back into the seam.
\item Lastly, the blockers, after cutting off the line should close shoulders to prevent the jammer cutting through. From this position they may trap the jammer.
\end{itemize}

Blockers may work on this drill using lateral or rotating motions.
\end{minipage}
